\section{Results}
\label{sec:results}


\begin{table*}[t]
\caption{
Full-duplex interaction and human evaluation on Full-Duplex-Bench v1.0.
RL-Seamless denotes \texttt{personaplex-rl-seamless}, and $c$ is the average number of interaction frames advanced per backbone wake.
Syn. and Can. denote the Synthetic and Candor pause sets.
TOR denotes turn-over rate; Freq., backchannel frequency; JSD, Jensen--Shannon divergence of backchannel timing; Lat., response latency; and Judge, the \textit{gpt-5.6-luna} response-appropriateness score for interruption responses (reasoning effort: none).
NMOS and DMOS denote naturalness and dialog mean opinion scores, reported as mean $\pm$ 95\% confidence interval.
}
\label{tab:full_duplex_results}
\centering
\setlength{\tabcolsep}{3.2pt}
\renewcommand{\arraystretch}{1.05}
\resizebox{\textwidth}{!}{%
\begin{tabular}{l
                c
                cc
                ccc
                cc
                ccc
                cc}
\toprule
& \multicolumn{1}{c}{Efficiency}
& \multicolumn{2}{c}{Pause}
& \multicolumn{3}{c}{Backchannel}
& \multicolumn{2}{c}{Turn-taking}
& \multicolumn{3}{c}{Interruption}
& \multicolumn{2}{c}{Human} \\
\cmidrule(lr){2-2}
\cmidrule(lr){3-4}
\cmidrule(lr){5-7}
\cmidrule(lr){8-9}
\cmidrule(lr){10-12}
\cmidrule(lr){13-14}
Model
& $c \uparrow$
& Syn. $\downarrow$
& Can. $\downarrow$
& TOR $\downarrow$
& Freq.
& JSD $\downarrow$
& TOR $\uparrow$
& Lat. $\downarrow$
& TOR $\uparrow$
& Judge $\uparrow$
& Lat. $\downarrow$
& NMOS $\uparrow$
& DMOS $\uparrow$ \\
\midrule
Moshi
& 1.00
& 0.993 & 1.000
& 0.582 & 0.184 & 0.679
& 1.000 & 0.000
& 0.765 & 3.359 & 1.229
& $3.52 \pm 0.15$ & $2.86 \pm 0.20$ \\

PersonaPlex
& 1.00
& 0.307 & 0.343
& 0.200 & 0.058 & 0.824
& 0.941 & 0.285
& 0.930 & 4.328 & 0.218
& $3.90 \pm 0.12$ & $4.10 \pm 0.13$ \\

RL-Seamless
& 1.00
& 0.299 & 0.301
& 0.218 & 0.163 & 0.740
& 1.000 & 0.124
& 0.995 & 4.286 & 0.281
& $3.93 \pm 0.12$ & $4.09 \pm 0.14$ \\
\midrule
DiffuPlex-\textsc{Listen}
& 1.91
& 0.277 & 0.343
& 0.109 & 0.140 & 0.753
& 1.000 & 0.173
& 0.990 & 4.328 & 0.329
& $3.92 \pm 0.11$ & $4.08 \pm 0.13$ \\

DiffuPlex-\textsc{Speak}
& 2.23
& 0.263 & 0.347
& 0.145 & 0.154 & 0.738
& 1.000 & 0.156
& 0.995 & 4.357 & 0.347
& $3.54 \pm 0.12$ & $4.07 \pm 0.13$ \\
\bottomrule
\end{tabular}%
}
\end{table*}
\subsection{Full-Duplex Interaction}
\label{subsec:full_duplex_results}

Table~\ref{tab:full_duplex_results} shows that DiffuPlex reduces backbone wake frequency while remaining close to the autoregressive backbone in full-duplex interaction behavior.
We use $c$ to denote the average number of interaction frames advanced per backbone wake; standard autoregressive generation has $c=1$.
DiffuPlex-\textsc{Listen} and DiffuPlex-\textsc{Speak} increase this to $c=1.91$ and $c=2.23$, respectively, while remaining close to \texttt{personaplex-rl-seamless} across pause handling, backchanneling, turn-taking, and user interruption.
Results on Full-Duplex-Bench v1.5 are provided in Appendix~\ref{asubsec:additional_full_duplex}.

The two settings expose a trade-off in how aggressively predicted futures are consumed.
DiffuPlex-\textsc{Listen} closely preserves both speech quality and conversational quality, with an NMOS of $3.92\pm0.11$ and a DMOS of $4.08\pm0.13$, compared with $3.93\pm0.12$ and $4.09\pm0.14$ for \texttt{personaplex-rl-seamless}.
DiffuPlex-\textsc{Speak} increases $c$ to 2.23 while maintaining a DMOS of
$4.07\pm0.13$, but its NMOS decreases to $3.54\pm0.12$, remaining close to Moshi at
$3.52\pm0.15$.
We analyze the source of this speech-quality degradation in Section~\ref{subsec:draft_analysis}.

Fewer backbone wakes do not necessarily imply lower runtime, because each DiffuPlex wake predicts over a wider future span and is therefore more expensive than a single autoregressive wake.
We next examine whether the reduction in wake frequency outweighs this increased per-wake cost.

\begin{table*}[t]
\caption{
Measured inference efficiency on the fixed 60-clip runtime subset.
Streaming wall is the elapsed wall-clock time of the deployment-path streaming loop, including CPU/Python control, synchronization, and GPU execution.
E2E GPU and Core LM denote accumulated GPU time for the full inference pipeline and the language-model backbone, respectively.
Spd. denotes speedup relative to RL-Seamless, and Rel. time/wake is the Core LM time per backbone wake normalized by RL-Seamless.
Wake-frame p99 is the 99th-percentile wall-clock latency of frames that invoke the backbone, and $>80$ms is the fraction of such frames exceeding the 80ms interaction interval.
}
\label{tab:runtime_results}
\centering
\setlength{\tabcolsep}{4pt}
\renewcommand{\arraystretch}{1.05}
\resizebox{\textwidth}{!}{%
\begin{tabular}{l
                c
                cc
                cc
                cc
                ccc}
\toprule
& \multicolumn{1}{c}{Wake}
& \multicolumn{2}{c}{Streaming wall}
& \multicolumn{2}{c}{E2E GPU}
& \multicolumn{2}{c}{Core LM}
& \multicolumn{3}{c}{Wake behavior} \\
\cmidrule(lr){2-2}
\cmidrule(lr){3-4}
\cmidrule(lr){5-6}
\cmidrule(lr){7-8}
\cmidrule(lr){9-11}
Model
& $c \uparrow$
& Time (s) $\downarrow$
& Spd. $\uparrow$
& Time (s) $\downarrow$
& Spd. $\uparrow$
& Time (s) $\downarrow$
& Spd. $\uparrow$
& Rel. time/wake
& p99 (ms) $\downarrow$
& $>80$ms $\downarrow$ \\
\midrule
Moshi
& 1.00
& 638.7 & 1.168$\times$
& 638.3 & 1.168$\times$
& 550.1 & 1.194$\times$
& 0.84$\times$
& -- & -- \\

PersonaPlex
& 1.00
& 744.1 & 1.002$\times$
& 743.7 & 1.002$\times$
& 656.8 & 1.000$\times$
& 1.00$\times$
& -- & -- \\

RL-Seamless
& 1.00
& 745.7 & 1.000$\times$
& 745.3 & 1.000$\times$
& 656.8 & 1.000$\times$
& 1.00$\times$
& 46.85 & 0\% \\
\midrule
DiffuPlex-\textsc{Listen}
& 1.98
& 512.4 & 1.455$\times$
& 511.1 & 1.458$\times$
& 408.4 & 1.608$\times$
& 1.23$\times$
& 61.34 & 0\% \\

DiffuPlex-\textsc{Speak}
& 2.26
& 469.1 & 1.590$\times$
& 467.8 & 1.593$\times$
& 365.5 & 1.797$\times$
& 1.26$\times$
& 61.41 & 0\% \\
\bottomrule
\end{tabular}
}
\end{table*}

\subsection{Runtime Efficiency}
\label{subsec:runtime_results}

Table~\ref{tab:runtime_results} shows that the reduction in backbone wakes translates into substantial measured speedups.
On the fixed 60-clip runtime subset, DiffuPlex-\textsc{Listen} and DiffuPlex-\textsc{Speak} advance $c=1.98$ and $c=2.26$ interaction frames per backbone wake, respectively.
This reduces accumulated Core-LM time from 656.8s for \texttt{personaplex-rl-seamless} to 408.4s and 365.5s, corresponding to $1.608\times$ and $1.797\times$ speedups.
The reduction also translates into deployment-path wall-clock speedups of $1.455\times$ and $1.590\times$ for DiffuPlex-\textsc{Listen} and DiffuPlex-\textsc{Speak}, respectively.
Considering device-side computation alone, this yields end-to-end GPU speedups of $1.458\times$ and $1.593\times$.
Moshi predicts only assistant-side tokens, whereas the PersonaPlex family jointly models both user and assistant streams, giving Moshi a lighter per-frame prediction workload.

These gains do not come from making each DiffuPlex backbone invocation cheaper.
The pooled measurements imply 1.23$\times$ and 1.26$\times$ as much Core LM time per
backbone wake for DiffuPlex-\textsc{Listen} and DiffuPlex-\textsc{Speak}, respectively.
Instead, masked-future prediction amortizes this larger per-wake cost over multiple interaction
frames, reducing accumulated Core LM time while sustaining the same 12.5Hz interaction stream.
Meeting the 80ms deadline establishes real-time feasibility, but systems that satisfy the same
deadline can still differ substantially in accumulated GPU time.
DiffuPlex therefore reduces accumulated GPU and Core LM time over the course of an interaction.

The larger individual wakes must still satisfy this real-time constraint.
The 99th-percentile wall-clock latency of frames that invoke the backbone is 61.34ms for DiffuPlex-\textsc{Listen} and 61.41ms for DiffuPlex-\textsc{Speak}, and no measured wake frame exceeds 80ms.
Thus, DiffuPlex reduces accumulated Core LM time while preserving the original real-time interaction cadence.

\subsection{Masked Future Prediction vs.\ Deferred Execution}
\label{subsec:deferred_results}

We next ask whether similar efficiency gains can be obtained without predicting future interaction, simply by skipping autoregressive backbone calls and emitting silence during skipped frames.
Figure~\ref{fig:deferred_tradeoff} compares DiffuPlex with the \textsc{PAD}, \textsc{Silero}, and \textsc{Smart Turn} deferred-execution baselines over $R\in\{2,4,6,8\}$.
Increasing $R$ allows up to $R-1$ frames to be skipped between backbone calls, improving wall-clock speedup but increasingly degrading interaction behavior.
In particular, deferred execution suppresses backchanneling and reduces assistant speaking activity, while turn-taking latency increases substantially for several controllers.
For example, at $R=6$, \textsc{PAD} reaches a $2.05\times$ wall-clock speedup, but backchannel frequency drops from 0.163 to 0.048 and turn-taking latency increases from 0.124s to 0.423s relative to the backbone.

DiffuPlex largely mitigates this trade-off by predicting interaction between backbone invocations rather than replacing skipped frames with silence.
DiffuPlex-\textsc{Listen} and DiffuPlex-\textsc{Speak} achieve $1.455\times$ and $1.590\times$ wall-clock speedups while remaining much closer to the original backchannel, turn-taking, and speaking behavior.
These results show that deferring backbone calls while emitting silence is insufficient in these baselines; predicting what occurs between backbone invocations better preserves full-duplex behavior.
DiffuPlex can further adjust this efficiency-interaction trade-off through its confidence and revision thresholds, which we analyze in Appendices~\ref{asubsec:confidence_retention} and~\ref{asubsec:revision}.


\begin{figure}[t]
    \centering
    \begin{subfigure}[t]{0.32\columnwidth}
        \centering
        \includegraphics[width=0.90\linewidth]{Figures/Figure3-1.pdf}
        \caption{Turn-taking latency ($\downarrow$)}
        \label{fig:deferred_turn}
    \end{subfigure}
    \hfill
    \begin{subfigure}[t]{0.32\columnwidth}
        \centering
        \includegraphics[width=0.90\linewidth]{Figures/Figure3-2.pdf}
        \caption{Backchannel frequency}
        \label{fig:deferred_bc}
    \end{subfigure}
    \hfill
    \begin{subfigure}[t]{0.32\columnwidth}
        \centering
        \includegraphics[width=0.90\linewidth]{Figures/Figure3-3.pdf}
        \caption{Assistant speaking fraction}
        \label{fig:deferred_speaking}
    \end{subfigure}
\caption{
Efficiency-interaction trade-offs under deferred execution.
Deferred baselines vary $R\in\{2,4,6,8\}$, and deferred frames are forced to silence.
Silero and Smart Turn replay precomputed controller decisions; controller inference cost is measured separately.
}
    \label{fig:deferred_tradeoff}
\end{figure}

\begin{table}[t]
\caption{
General spoken-language capability. Higher is better.
}
\label{tab:general_results}
\centering
\small
\setlength{\tabcolsep}{4.5pt}
\renewcommand{\arraystretch}{1.05}
\begin{tabular}{lcccc}
\toprule
& \multicolumn{2}{c}{VoiceBench}
& \multicolumn{2}{c}{OpenAudioBench} \\
\cmidrule(lr){2-3}
\cmidrule(lr){4-5}
Model
& $s{\rightarrow}t$
& $s{\rightarrow}s$
& $s{\rightarrow}t$
& $s{\rightarrow}s$ \\
\midrule
\texttt{personaplex-rl-seamless}
& 31.95 & 31.61
& 35.66 & 34.03 \\

DiffuPlex-\textsc{Listen}
& 31.93 & 31.52
& 36.71 & 33.53 \\

DiffuPlex-\textsc{Speak}
& 31.20 & 31.04
& 35.82 & 32.12 \\
\bottomrule
\end{tabular}
\end{table}
\begin{table}[t]
\caption{
Ablation of confidence-based retention and draft revision on Full-Duplex-Bench v1.0.
}
\label{tab:inference_ablation}
\centering
\small
\setlength{\tabcolsep}{3.5pt}
\renewcommand{\arraystretch}{1.03}
\begin{tabular}{llccccc}
\toprule
Policy & Variant
& $c\uparrow$
& BC Freq.
& Turn Lat.$\downarrow$
& Int. Judge$\uparrow$
& Int. Lat.$\downarrow$ \\
\midrule
\multirow{3}{*}{\textsc{Listen}}
& DiffuPlex
& 1.91 & .140 & .173 & 4.328 & .329 \\
& w/o confidence
& 2.11 & .191 & .170 & 4.398 & .314 \\
& w/o revision
& 2.49 & .078 & .440 & 4.372 & .744 \\
\midrule
\multirow{3}{*}{\textsc{Speak}}
& DiffuPlex
& 2.23 & .154 & .156 & 4.357 & .347 \\
& w/o confidence
& 2.80 & .191 & .209 & 3.550 & .281 \\
& w/o revision
& 3.08 & .080 & .480 & 4.251 & .515 \\
\bottomrule
\end{tabular}
\end{table}

\subsection{General Spoken-Language Capability}
\label{subsec:general_results}

Table~\ref{tab:general_results} shows that DiffuPlex largely preserves the general spoken-language capability of \texttt{personaplex-rl-seamless}.
DiffuPlex-\textsc{Listen} remains nearly unchanged on
VoiceBench and OpenAudioBench, while DiffuPlex-\textsc{Speak} shows moderate degradation,
most notably on OpenAudioBench $s\!\to\!s$.
Full per-task results are provided in Appendix~\ref{asubsec:additional_general}.


\subsection{Analysis of Future Drafts}
\label{subsec:draft_analysis}


\textbf{Rolling draft behavior.}
We first examine how confidence filtering and revision affect the use of predicted futures.
As shown in Table~\ref{tab:inference_ablation}, the policy-specific no-confidence variants allow more aggressive
retention, with different effects across the two policies.
For \textsc{Speak}, $c$ increases from 2.23 to 2.80, but the interruption Judge score decreases from 4.357 to 3.550.
For \textsc{Listen}, where retention remains restricted to PAD states, removing confidence increases $c$ more modestly from 1.91 to 2.11 while leaving interaction behavior similar.
Removing revision increases $c$ from 1.91 to 2.49 for \textsc{Listen} and from 2.23 to 3.08 for \textsc{Speak}, but substantially reduces responsiveness: backchannel frequency falls to 0.078 and 0.080, while turn-taking latency increases to 0.440s and 0.480s, respectively.
Thus, the retention constraints limit aggressive future consumption, particularly for \textsc{Speak},
while revision is important for preserving responsiveness to newly observed user speech.
Sensitivity to the confidence threshold $\tau_{\rm conf}$ and revision threshold $\tau_{\rm rev}$ is reported in Appendices~\ref{asubsec:confidence_retention} and~\ref{asubsec:revision}.

We next examine why DiffuPlex wakes the backbone again after producing a draft.
A wake occurs either when the retained draft is exhausted or when new user speech triggers revision.
When the draft is exhausted, DiffuPlex advances only 1.17 frames on average for \textsc{Listen} and 1.48 for \textsc{Speak}; 98.2\% and 81.5\% of these wakes, respectively, occur after a single frame.
Revision-triggered wakes instead occur after several draft frames have already been consumed, with mean advancements of 4.24 and 3.79 frames and medians of 4 and 3.
In other words, a short wake interval usually means that only a short prefix was retained in the first place, rather than that a longer retained draft was quickly revised.
When revision does occur, several retained frames have typically been consumed.
Additional wake-level diagnostics and revision analyses are provided in
Appendix~\ref{asubsec:wake_dynamics}.

\textbf{Why does future speech degrade?}
We next investigate the speech-quality degradation of DiffuPlex-\textsc{Speak}.
\texttt{personaplex-rl-seamless} models assistant audio using a semantic codebook, cb0, followed by seven acoustic codebooks, cb1--7.
Because cb1--7 are natively delayed by one temporal step, the cb0 corresponding to a given physical audio frame appears one temporal position earlier in the model sequence and is therefore already available when the autoregressive model predicts cb1--7 for that frame.
DiffuPlex preserves this native delay pattern, but at future depths $h\geq2$ all assistant streams are masked, so future cb1--7 must be predicted without the same-frame cb0 that would normally condition them.
To test whether this missing conditioning contributes to the degradation, we use the 199 valid user-interruption examples from Full-Duplex-Bench v1.0 and replace only the masked cb0 with the ground-truth cb0 for the physical frame whose cb1--7 are being predicted.
During assistant speech, this reduces the cb1--7 negative log-likelihood from 5.702 to 5.611
($\Delta=0.091$, 95\% CI [0.081, 0.101]), compared with a much smaller reduction during silence ($\Delta=0.029$).
These results show that same-frame semantic conditioning has a substantially larger effect on acoustic-token prediction during speech than during silence.

We then test whether restoring autoregressive audio-token conditioning improves generated audio.
Starting from each DiffuPlex-\textsc{Speak} output, we teacher-force its generated text and cb0 while recomputing the sequence frame by frame through the original autoregressive backbone.
We then use the resulting temporal representations to regenerate only cb1--7 with the unchanged depth transformer.
UTMOSv2, a neural predictor of perceptual speech quality on a 5-scale~\citep{10832315}, increases from 2.646 to 2.894.
Regenerating cb0--7 while keeping only the generated text fixed yields 2.903, compared with 2.919 for fully autoregressive generation.
Thus, most of the measured quality loss is recovered by restoring autoregressive audio-token conditioning even when the text and cb0 generated by DiffuPlex are unchanged.
Additional intervention controls are in
Appendix~\ref{asubsec:future_speech}.
